
\begin{proposition}[Exact pruning and compact reinsertion classify every reduced all-Hopf support]\label{cdp:st:C.2.UN.60}
Let $S$ be a specified compact connected reduced SNC surface whose
smooth components are primary or secondary Hopf surfaces. Its
elliptic gluing maps are part of the specified surface. The following
finite geometric test is necessary and sufficient for $S$ to occur
as the full central fibre of a proper degeneration of connected
complex two-tori with smooth total space.

\begin{enumerate}
\item The actual normal lines on every double curve are dual, and
the incidence graph has exactly one independent cycle.
\item If a diagonal algebraic-dimension-zero component occurs, the
whole support is its diagonal cycle and its integral lattice product
$M$ is the identity, as in \ref{cdp:st:C.2.UN.46}.
\item In the remaining case every algebraic-dimension-zero vertex
is a nonlinear leaf. Every other vertex has its actual elliptic
reduction to $\PP^1$, and every multiple fibre of that reduction
is one of its marked double curves.
\item Let $m_e$ be the exact order of either normal line at an edge,
and let $h_v$ be zero on the cycle and the sum of the $m_e$ along
the path from the cycle to $v$. There is an integer $R>0$ such that
every nonlinear leaf has $h_v=R$ and every other vertex has $h_v<R$.
If there is no nonlinear leaf, any $R>\max h_v$ passes this test.
\item Perform the exact leaf-pruning operations specified below,
including the full integer secondary framing and all torsor maps.
The resulting marked Hopf/ruled core cycle has integral
complement-lattice product $M=I$.
\end{enumerate}

In the second case the tests referring to elliptic tree vertices are
unneeded. In the remaining case the pruning operations always exist
under the preceding geometric requirements. Each terminal step is
one of two explicit operations: the inverse logarithmic chart of
\ref{cdp:st:C.2.UN.50} for a primary or secondary Hopf pair; or removal of the
product pair $(F\times\PP^1,F\times\{\infty\})$. The latter leaves
the parent surface unchanged. After its children have been removed,
every vertex has exactly one of these two forms. The pruned core
and its test are independent, up to the stated actual marked
isomorphisms, of the order of pruning disjoint branches.

Every support passing the test has a realization over $\PP^1$
with smooth compact total threefold. It retains all the original
components, their analytic torsors and every specified elliptic
gluing map. Its other singular support can be taken to be the
complex conjugate of its pruned core cycle. All reduced fibre
supports have Euler characteristic zero.
\end{proposition}

\begin{proof}
\par\medskip\noindent\textbf{The exact one-point torsor and its marked compact surface.}
We first prove the classification used at a terminal pruning step.
Let $q:T\to\PP^1$ be a fixed-Jacobian, identity-monodromy elliptic
quasi-bundle, with just one possible multiple fibre $mJ$, marked
over infinity. Its actual normal-root charts make it a principal
$F$-torsor on
\[
                 \mathcal B=\PP^1[(\infty,m)].          \tag{C.PR.1}\label{cdp:eq:C.PR.1}
\]
Choose an oriented period basis $\lambda_1,\lambda_2$ of
$F=\C/\Lambda_F$. The exact sheaf sequence used here is
\[
 \begin{gathered}
 0\longrightarrow\OO_{\mathcal B}
 \xrightarrow{\ a\mapsto
 (e^{2\pi i\lambda_2a},e^{-2\pi i\lambda_1a})\ }
 (\OO_{\mathcal B}^{*})^2
 \xrightarrow{\ \Phi\ }F_{\mathcal B}\longrightarrow0,\\
 \Phi(u,v)=
 \frac{\lambda_1\log u+\lambda_2\log v}{2\pi i}
                          \pmod{\Lambda_F}.            \tag{C.PR.2}\label{cdp:eq:C.PR.2}
 \end{gathered}
\]
Changing a logarithm changes the numerator by a lattice period.
Conversely a pair in the kernel can be lifted to logarithms,
adjusted by integral periods, and then lies on the displayed
complex line. The intersection of that line with the integral
two-period lattice is zero because
$\lambda_2/\lambda_1$ is not real. Thus the additive kernel and
every map in \eqref{cdp:eq:C.PR.2} are exact.

Finite stabilizer invariants are exact in characteristic zero, so
$H^j(\mathcal B,\OO)=H^j(\PP^1,\OO)=0$ for $j=1,2$.
Moreover
\[
 \operatorname{Pic}(\mathcal B)=\Z[\mathcal R],
                 \qquad \mathcal R^m=\OO_{\PP^1}(\infty).
                                                        \tag{C.PR.3}\label{cdp:eq:C.PR.3}
\]
Indeed every line has a unique stabilizer weight modulo $m$;
subtract that power of the canonical root line and it descends
to $\PP^1$. The displayed relation then gives the asserted
infinite cyclic group, with its canonical sections retained.
The long exact sequence of \eqref{cdp:eq:C.PR.2} therefore identifies the
actual $F$-torsor class with a unique vector in
$\operatorname{Pic}(\mathcal B)^2=\Z^2$.

If this vector is zero, the torsor is the product
$F\times\mathcal B$. The inertia must act freely on its total
space because the actual connected root cover of $mJ$ is
unramified over the original surface tube. A nontrivial inertia
would act trivially on its product fibre. Thus $m=1$ and
$T=F\times\PP^1$, with $J=F\times\{\infty\}$.

For a nonzero vector, choose its primitive direction and an
oriented period basis so that it is $(-n,0)$, $n>0$. This is
an integral change of the actual fibre basis, not removal of
its direction. The torsor is the pushout of
$(\mathcal R^{-n},\OO)$ by \eqref{cdp:eq:C.PR.2}. Gauge the second nonzero
line coordinate to one. The remaining ambiguity is multiplication
of the first coordinate by
$\gamma^{\Z}$, where
$\gamma=\exp(2\pi i\lambda_2/\lambda_1)$ and $|\gamma|<1$.
Consequently the actual total surface is
\[
                         (\mathcal R^{-n})^\times/\gamma^{\Z}.
                                                        \tag{C.PR.4}\label{cdp:eq:C.PR.4}
\]
Write $\mathcal B=\PP(m,1)$ with coordinates $(Z,W)$ of
weights $(m,1)$ and coarse coordinate $U=Z/W^m$.
The nonzero total line of $\mathcal R^{-n}$ is
\[
 [((\C^2\setminus0)\times\C^*_\ell)/
              \lambda\cdot(Z,W,\ell)
                =(\lambda^mZ,\lambda W,\lambda^{-n}\ell)].
                                                        \tag{C.PR.5}\label{cdp:eq:C.PR.5}
\]
Setting $\ell=1$ leaves the finite group $\mu_n$ with weights
$(m,1)$. A scalar multiplication by $\gamma$ lifts, after choosing
$\beta^n=\gamma$, to $(Z,W)\mapsto(\beta^mZ,\beta W)$.
Thus \eqref{cdp:eq:C.PR.4} is exactly
$H_m^{\rm diag}(\beta)/\mu_n$, with the marked axis $W=0$.
Freeness of the original inertia is exactly
$\gcd(m,n)=1$. Changing the chosen root of $\gamma$ changes
the infinite generator by the already included finite group.
Every primitive fibre basis, the full integer $n$, and the marked
axis identification have been retained.

This also explains why reduction of $n$ modulo $m$ is impermissible:
the entire root-line degree in \eqref{cdp:eq:C.PR.4}, rather than only its
stabilizer character, is the torsor class.

\par\medskip\noindent\textbf{The actual inverse operation at a Hopf leaf.}
For a nonlinear leaf use its full normal form from \ref{cdp:st:C.2.UN.43};
for an elliptic terminal surface use \eqref{cdp:eq:C.PR.4}. Both have the
same marked normal-line data. Put $r=m$, choose $k,l$ with
$mk-nl=1$, and write $\gamma=\beta^n$. On the parent the
actual coarse root chart and its inverse product chart are
\[
 \begin{gathered}
 (u,X)\sim(\gamma^m u,\gamma^lX),\qquad p=X^m u^{-l},\\
 W=X^{-n}u^k,\qquad
                u=p^nW^m,\quad X=p^kW^l.
 \end{gathered}                                        \tag{C.PR.6}\label{cdp:eq:C.PR.6}
\]
These are precisely (C.SN.3)--(C.SN.5), with the given dual-normal
map and elliptic gluing used to choose the same connected root
cover. Replace this marked multiple-fibre tube by
$(p,W)\in\Delta\times(\C^*/\gamma^{\Z})$ through the displayed
isomorphism on $p\ne0$. This is the actual inverse logarithmic
transform. Its base point becomes an ordinary elliptic fibre.
Every other chart and transition of the parent is unchanged.
In particular the parent remains the same fixed-Jacobian
identity-monodromy quasi-bundle off this one point, with its full
analytic torsor retained by the displayed transition map.

The choices do not discard gluing data. A change of a root
trivialization or a normal-product constant changes \eqref{cdp:eq:C.PR.6} by
a translation gauge extending over the inserted product disc.
The primitive nonzero direction in \eqref{cdp:eq:C.PR.3} fixes the angular
direction; the other period can change by an integral multiple,
which changes the generator by its same finite deck group.
Translations of the marked axis are retained in the transition.
This is the exact choice comparison already proved for the
full inverse torsor in \ref{cdp:st:C.2.UN.49} and \ref{cdp:st:C.2.UN.50}.

Remove terminal vertices inward from every attached tree.
All original nonterminal vertices have the elliptic reduction
of \ref{cdp:st:C.2.UN.58}. Removing a child removes exactly its marked
multiple fibre by \eqref{cdp:eq:C.PR.6}, or leaves an ordinary fibre unchanged
in the product-cap case. Hence, after all children of a vertex
are removed, its only possible multiple fibre is its parent
mark. The one-point classification above applies to that
actual modified surface. This proves that every terminal step
has one of the two asserted forms.

Operations in distinct discs of one base keep the same transition
atlas on their complement and hence commute. Operations in disjoint
subtrees do likewise. Every allowable pruning order differs by
interchanging such operations; a parent is processed only after
its children. This proves independence through actual marked
torsor isomorphisms, not merely through the same degree vector.
Keep the entire maps \eqref{cdp:eq:C.PR.6}, the primitive periods, the integer
$n$ and all remaining marked points as the record of each step.

\par\medskip\noindent\textbf{The exact two-point core models.}
At the end each cycle vertex is still a connected fixed-$F$
quasi-bundle over $\PP^1$, with only its two marked possible
multiple fibres, of orders $a,b$. Put $g=\gcd(a,b)$,
$A=a/g$, $B=b/g$. There is the actual representable finite
etale orbicurve map
\[
 \PP^1[(0,A),(\infty,B)]
       \longrightarrow \PP^1[(0,a),(\infty,b)],
                         \qquad t=z^g .                \tag{C.PR.7}\label{cdp:eq:C.PR.7}
\]
At zero its root uniformizers have map $v=u$ and inertia
$\mu_A\subset\mu_a$; the infinity chart has the analogous
inclusion. Thus it is etale as a map of these root curves.
The source, since $A,B$ are coprime, is the weighted projective
line $\PP(B,A)$, with coarse coordinate
$z=X^A/Y^B$. Its Picard group is generated by $\OO(1)$:
the two root divisors have classes $B$ and $A$, and Bezout
expresses their generator. Its higher coherent groups vanish.

Pull back the actual $F$-torsor. The same sequence \eqref{cdp:eq:C.PR.2}
therefore describes it by an integral vector. A nonzero vector
$(-n,0)$ gives the explicit nonzero line $\OO(-n)$, and the
gauge \eqref{cdp:eq:C.PR.5}, now of weights $(B,A)$, gives the actual diagonal
Hopf cover with contraction $(\beta^BX,\beta^AY)$ and finite
action $\mu_n(B,A)$. The inertia is free precisely when
$\gcd(n,A)=\gcd(n,B)=1$. Its cyclic descent preserves both
marked axes. Lifts to a central primary cover commute with a
contraction and preserve each axis, so they have the form
$(Xf(X,Y),Yh(X,Y))$; contraction equivariance forces the
nonvanishing units $f,h$ to be constant. Thus the full quotient
is the actual marked diagonal primary or secondary Hopf surface.

For the zero vector, free inertia forces $A=B=1$ and the cover
is $F\times\PP^1$. Its cyclic descent is
$(y,z)\mapsto(y+\tau,\zeta_g z)$: a holomorphic map
$\PP^1\to F$ is constant, and the principal-$F$ action excludes
a linear part. Freeness at $0,\infty$ makes $\tau$ have exact
order $g$. The quotient is exactly
\[
        \PP(\OO_E\oplus L),\quad
        E=F/\langle\tau\rangle,\quad
        L=(F\times\C)/((y,v)\sim(y+\tau,\zeta_gv)).
                                                        \tag{C.PR.8}\label{cdp:eq:C.PR.8}
\]
This exhibits its two marked sections, the exact torsion order
of $L$, the full fibre map $q=z^g$ on the cover, and both
elliptic identifications. It also gives the actual finite
projective-bundle map from $F\times\PP^1$.
Thus the pruned core is precisely one of the marked diagonal
Hopf/ruled cycles of \ref{cdp:st:C.2.UN.57}, with all its data retained.

\par\medskip\noindent\textbf{Necessity of the core lattice test through actual group maps.}
For an actual torus degeneration, the graph, normal and height
requirements are \ref{cdp:st:C.2.UN.54} and \ref{cdp:st:C.2.UN.58}. It remains to compare
the nearby group to that of the pruned core; no contraction of
the given holomorphic total space is required.

The complement of the marked axis in a terminal diagonal Hopf
pair is exactly $\C\times F$. For a nonlinear pair the explicit
proper smooth-pair deformation of \ref{cdp:st:C.2.UN.49} identifies its compact
core and boundary with the diagonal pair. The product ruled cap
has the same complement $F\times\C$. In the primitive boundary
basis of \ref{cdp:st:C.2.UN.50} the actual homomorphism of this compact
piece is
\[
        \Z^3_{\delta,e,h}\longrightarrow\Z^2_{e,h},
                  \delta\longmapsto0,\quad
                  e\longmapsto e,\quad h\longmapsto h. \tag{C.PR.9}\label{cdp:eq:C.PR.9}
\]
For the Hopf case the exact matrix
$[(-l,k,0),(0,0,1),(m,-n,0)]$ has determinant one;
the normal reversal in $xy=t$ turns the killed Hopf meridian
into precisely $\delta$. The inverse product filling
$F\times D^2$ has the identical homomorphism \eqref{cdp:eq:C.PR.9},
with the same two retained generators. For a product cap this
is the literal product-disc boundary map.

Van Kampen applied to the unchanged exterior and this boundary
gives two pushout diagrams with identical maps. Their resulting
groups are canonically isomorphic by the identity on the exterior
and the two displayed fibre generators. Repeat the comparison,
retaining every remaining collar. The actual tubular comparison
of \ref{cdp:st:C.2.UN.49} carries these maps to an arbitrary original smoothing.
Consequently the nearby group of the core plumbing is the
original $\Z^4$. The exact core calculation of \ref{cdp:st:C.2.UN.57} gives
$\Lambda\rtimes_M\Z$, with $\Lambda$ injected. Its abelianness
forces $M=I$. This proves necessity using the actual maps needed
for the test.

\par\medskip\noindent\textbf{The compact core family with a prescribed remaining period order.}
For sufficiency first handle the all-algebraic-dimension-zero
diagonal branch by \ref{cdp:st:C.2.UN.48}. Otherwise prune as above and assume
the core test passes. Its components have retained elliptic
quasi-bundles. In the common complement
$G=\C^2/\Lambda$ of \ref{cdp:st:C.2.UN.57} their regular fibres are translates
of the compact elliptic subgroup
\[
       F=(W\cap iW)/(\Lambda\cap W\cap iW),
                     \qquad W=\operatorname{span}_{\R}\Lambda.
                                                        \tag{C.PR.10}\label{cdp:eq:C.PR.10}
\]
Their lattice is primitive of rank two: their actual quotient
in the open two-boundary component is $\C^*$.
The root-cover maps already recorded identify every later
inserted regular fibre with this same elliptic group.

In \ref{cdp:st:C.2.UN.57} replace $V_z(w)=w+zJ\overline w$ by
$V_{t^R}(w)=w+t^RJ\overline w$, retaining the term
$d\log t$ and the normal neck equation $xy=t$.
The exact real determinant becomes
\[
 C(1+|t|^{2R})
 \left(1+|t|^{2R}
       +\frac{d}{2\pi}(|t|^{2R}-1)\log|t|\right)>0.
                                                        \tag{C.PR.11}\label{cdp:eq:C.PR.11}
\]
The product $(|t|^{2R}-1)\log|t|$ is nonnegative.
Thus the same actual period parallelepipeds and both endpoint
plumbings give a smooth compact torus family on $\PP^1$ with
the prescribed core at zero and its conjugate at infinity.
At infinity divide the frame by $t^R$; the displayed frame is
$J\overline w+s^Rw$, and the logarithmic term still has exactly
the sign and coefficient for the reduced neck $xy=s$.

Let $\lambda,\mu=c\lambda$ be an oriented primitive period basis
of the compact group in \eqref{cdp:eq:C.PR.10}, with $\operatorname{Im}c>0$.
The ordinary determinant in the actual vector frame is
\[
 \det_{\C}(V_{t^R}\lambda,V_{t^R}\mu)
       =2i\operatorname{Im}(c)\,\|\lambda\|^2 t^R\ne0
                      \quad(t\ne0).                    \tag{C.PR.12}\label{cdp:eq:C.PR.12}
\]
Here $\det(\lambda,J\overline\lambda)=-\|\lambda\|^2$.
This is the exact common core period, including its integral
class, rather than a chosen formal order.

\par\medskip\noindent\textbf{The recursive local model and the exact order drops.}
The actual logarithmic charts near a regular compact fibre give
a neighbourhood in the following rank-two period quotient:
\[
 \mathcal G=
 (\C_p\times\C_w^*\times\Delta_t)/
 \langle(p,w,t)\mapsto(p+\delta(t),\gamma(t)w,t)\rangle,
 \quad
 \delta(t)=t^\ell d(t),\quad d(0)\ne0,\quad|\gamma(0)|<1 .
                                                        \tag{C.PR.13}\label{cdp:eq:C.PR.13}
\]
Indeed choose the primitive first period as $(0,2\pi i)$ in
holomorphic linear coordinates; the second is then
$(\delta(t),\log\gamma(t))$. Its first coordinate is a unit
times their determinant. The charts of the core plumbing select
a bounded transverse disc in this cover and identify it with
the actual neighbourhood of the chosen central elliptic fibre.
Initially \eqref{cdp:eq:C.PR.12} gives $\ell=R$.

Every prescribed new primitive fibre basis is allowed. If
$\lambda'=a\lambda+b\mu$, $\mu'=c\lambda+d\mu$ with
$ad-bc=1$, the coordinate
$p'=p-b\delta V/(2\pi ia+b\log\gamma)$ and the rescaled
logarithmic coordinate
$V'=2\pi iV/(2\pi ia+b\log\gamma)$ make $\lambda'=(0,2\pi i)$.
The second displacement is
$2\pi i\,\delta/(2\pi ia+b\log\gamma)$.
The denominator is a holomorphic unit near zero.
Thus its order is unchanged and the full new fibre basis is
retained by an actual linear period map.

For a recorded Hopf step of weight $m$ and secondary integer $n$,
perform the weight-$(m,1)$ blowup of the base coordinates $(p,t)$
and tensor its nonzero line by $\OO(nE)$. In the free coarse
chart of \ref{cdp:st:C.2.UN.50} write $p=x^mZ$, $t=xW$.
A chosen holomorphic root $\beta(t)^n=\gamma(t)$ gives the exact
infinite map
\[
 (Z,W,x)\longmapsto
 \left(\beta(t)^m
    \bigl(Z+d(t)x^{\ell-m}W^\ell\bigr),
       \beta(t)W,\beta(t)^{-1}x\right).
                                                        \tag{C.PR.14}\label{cdp:eq:C.PR.14}
\]
The finite action has weights $(m,1,-1)$ modulo $n$ and is free,
because $\gcd(m,n)=1$. The added monomial has weight $m$,
so these maps commute. The inequality $\ell\ge m$ follows
from the height test. At $\ell=m$ the exceptional component
is exactly the recorded nonlinear Hopf surface after a constant
rescaling of $Z$. At $\ell>m$ it is exactly the recorded
diagonal pair. On the new child's complement of its parent curve
put $U=Z/W^m$. The actual deck map becomes
\[
        (U,W^n)\longmapsto
       \left(U+\frac{\delta(t)}{t^m},\,\gamma(t)W^n\right).
                                                        \tag{C.PR.15}\label{cdp:eq:C.PR.15}
\]
It has remaining displacement order $\ell-m$.
At any other marked point use $U-U_0$ as its base coordinate.
This retains the original marked point and every later child
transition. No extension of an arbitrary original smoothing's
elliptic quotient is being asserted.

For a recorded product cap, $m=1$ and $\ell\ge2$, because this
vertex was an elliptic vertex of height strictly smaller than $R$.
Blow up the actual ideal $(p,t)$ along its compact elliptic
curve in \eqref{cdp:eq:C.PR.13}. The deck map preserves this ideal.
On its exceptional normal directions it is the identity,
since $\delta'(0)=0$; hence the exceptional is precisely
$F\times\PP^1$ with its parent section. The chart $p=tU$
has the exact map
\[
                    (U,w)\longmapsto
          (U+\delta(t)/t,\gamma(t)w).                   \tag{C.PR.16}\label{cdp:eq:C.PR.16}
\]
Thus this step has exactly the same order drop one. The strict
parent surface is unchanged, since its centre is a Cartier
curve on that surface. The two other blowup coordinates give
the reduced normal-crossing equation $t=xy$.

Inductively, the model at a vertex of height $h$ has remaining
order $R-h$. All vertices with descendants are elliptic and
have positive remaining order. Only a nonlinear terminal step
has order zero after insertion. The height test is therefore
precisely what makes every recorded step in \eqref{cdp:eq:C.PR.14}--\eqref{cdp:eq:C.PR.16}
holomorphic with the required central type.

\par\medskip\noindent\textbf{Local properness and simultaneous compact replacement.}
We give the needed bounds also for the affine, noncompact base
of \eqref{cdp:eq:C.PR.13}. On its weight-$(m,1)$ blowup use the actual
canonical-section norm
\[
                  \rho(p,t)=
                  (|p|^2+|t|^{2m})^{1/(2m)}.           \tag{C.PR.17}\label{cdp:eq:C.PR.17}
\]
In the chart $p=x^mZ$, $t=xW$ this is
$|x|(|Z|^2+|W|^{2m})^{1/(2m)}$, the nonvanishing smooth
metric factor on the weighted exceptional charts. Translation
by $j\delta(t)$ gives two-sided norm distortion at most
$C(1+|j|)^{1/m}$, because $|\delta(t)|\le C|t|^m$.
For $\OO(nE)$ the exponent is $n/m$. On a sufficiently small
closed disc the scalar $|\gamma(t)|$ is uniformly smaller
than one. Its exponential contraction dominates this polynomial.
The infinite action is free and properly discontinuous, and
a high fixed iterate gives a norm-annulus representative for
every orbit, exactly as in \ref{cdp:st:C.2.UN.42}.

The old line coordinate is the new line divided by the $n$th
power of the canonical exceptional section. If an old image
stays in a compact old annular chart, putting both line norms
in annuli changes the base by at most a constant times
$|\delta(t)|(1+|\log\rho|)$. For $\rho<1$ use $\rho\ge|t|$;
this quantity is bounded and tends to zero as $t\to0$.
For $\rho\ge1$, a new base representative consequently obeys
\[
 |p_{\rm new}|\le B+C|\delta(t)|\log(2+|p_{\rm new}|).
                                                        \tag{C.PR.18}\label{cdp:eq:C.PR.18}
\]
It is uniformly bounded: a linear function eventually dominates
the logarithm. Thus the new annulus over the required base set
is compact, including the compact exceptional Hopf component.
This proves the local compactness comparison without assuming
that the affine base itself is compact.

Off the new Hopf component the canonical exceptional section
gives the actual isomorphism with the old model minus its
marked central elliptic curve. In a compact Hopf fundamental
domain, $x\to0$ gives $p=x^mZ$, $t=xW$ and old line coordinate
proportional to $x^{-n}$. Normalizing that coordinate takes
$O(1+|\log|x||)$ iterates. Every old base cluster point is zero,
because
$|x|^m+|t|^\ell(1+|\log|x||)\to0$.
The normal-annulus coordinate has a compact limit on the old
elliptic curve. Hence every old cluster point lies on that
curve and on no other point of the old model.

Choose disjoint small neighbourhoods of the finitely many
curves at each insertion stage. The new open piece is the
inserted compact component together with the inverse image
of the punctured old neighbourhood. The preceding cluster
comparison proves it is open and that all gluing graphs are
relatively closed. The resulting complex space is Hausdorff.
The bounded-annulus comparison over smaller closed
neighbourhoods and the unchanged compact exterior prove
properness. Product-cap steps are ordinary proper blowups.
Finite iteration therefore gives a smooth proper family over
a disc, glued by actual biholomorphisms to the unchanged
compact core family away from zero.

\par\medskip\noindent\textbf{Restoration of the prescribed support.}
On the strict parent of every Hopf insertion the chart is
exactly \eqref{cdp:eq:C.PR.6}, in the opposite direction: the old coordinate
is $W=x^{-n}$ and the full secondary integer is unchanged.
It restores the complete recorded elliptic torsor, not just
its rational Chern class. All untouched charts and all
other marked fibres remain their recorded charts. On a cap
step the actual product exceptional pair and its parent
section are the recorded ones. Fibre-basis maps retain
the linear parts of every elliptic gluing. The constant
translations are retained in the recorded transition maps;
equivalently the automorphisms
$(Z,W)\mapsto(a^mZ,aW)$ of either Hopf contraction realize
every required marked-axis translation. They commute with
the finite action. Product caps have the corresponding
elliptic translations and projective automorphisms fixing
the parent section.

Reverse the complete pruning record. At each step these
actual maps recover exactly the surface and marked gluing
removed at that step. Induction therefore restores the
specified $S$ in its full analytic isomorphism class.
The normal equations are $t=xW$ in the free weighted
coarse charts and $t=xy$ for an ordinary blowup, so every
component has full multiplicity one. The other fibres
are unchanged from the compact period family, and the
filling at infinity is its conjugate pruned core.
All these reduced supports have Euler zero: their
components are Hopf or elliptic ruled surfaces, their
double curves are elliptic, and there are no triple points.
This proves sufficiency and the compact assertion.

The exact comparisons are reproduced in
\path{code/E_hopf_tree_pruning.py} and its saved output:
primitive integral fibre-basis changes, both-point root Picard
lattices, the complete weighted deck and finite-action maps,
the nested remaining-order identities, and the product-cap
blowup. These checks retain the full integers in the pruning
record; the proof establishes them for arbitrary finite trees.
\end{proof}


